\documentclass[12pt,a4paper]{report}

% ============================================================
% PACKAGES
% ============================================================
\usepackage[T5]{fontenc}
\usepackage[utf8]{inputenc}
\usepackage[vietnamese]{babel}
\usepackage{lmodern}
\usepackage{microtype}

\usepackage{geometry}
\geometry{top=3cm, bottom=3cm, left=3.5cm, right=2.5cm}
\usepackage{setspace}
\onehalfspacing
\clubpenalty=10000
\widowpenalty=10000

\usepackage{amsmath, amssymb, amsthm, mathtools}
\usepackage{booktabs, adjustbox, multirow, array}
\usepackage{graphicx, caption, subcaption}
\captionsetup{font=small, labelfont=bf}

\usepackage[colorlinks=true, linkcolor=black,
citecolor=black, urlcolor=blue]{hyperref}
\usepackage{indentfirst}
\usepackage{enumitem}
\usepackage{tocbibind}

\usepackage{tikz}
\usepackage{pgfplots}
\pgfplotsset{compat=1.18}
\usetikzlibrary{shapes.geometric, arrows.meta, positioning,
	fit, backgrounds, calc, decorations.pathreplacing}
\tikzset{every picture/.style={line width=0.8pt}}

% ============================================================
% THEOREM ENVIRONMENTS
% ============================================================
\theoremstyle{plain}
\newtheorem{dinhly}{Định lý}[chapter]
\newtheorem{menh}{Mệnh đề}[chapter]

\theoremstyle{definition}
\newtheorem{dinhnghia}{Định nghĩa}[chapter]
\newtheorem{vidu}{Ví dụ}[chapter]

% ============================================================
% MATH OPERATORS
% ============================================================
\DeclareMathOperator*{\argmax}{argmax}
\DeclareMathOperator*{\argmin}{argmin}

% ============================================================
% TITLE PAGE
% ============================================================
\begin{document}
	\pagenumbering{roman}
	\pagestyle{plain}
	
	\begin{titlepage}
		\centering
		\vspace*{1cm}
		\textbf{TRƯỜNG ĐẠI HỌC KHOA HỌC}\\[0.2cm]
		\textbf{ĐẠI HỌC HUẾ}\\[0.8cm]
		\textbf{\Large KHOA CÔNG NGHỆ THÔNG TIN}\\[1.2cm]
		
		\includegraphics[width=2.5cm, keepaspectratio]{logo.png}\\[0.5cm]
		{\small\textit{Trường Đại Học Khoa Học -- Đại Học Huế}}\\[1.5cm]
		
		\textbf{\LARGE KHOÁ LUẬN TỐT NGHIỆP}\\[0.8cm]
		
		\textbf{\Large Chiến lược Phân đoạn Văn bản (Chunking)\\
			trong Hệ thống Truy xuất Tăng cường (RAG):\\
			Từ Lý thuyết đến Ứng dụng Tư vấn Tuyển sinh}\\[2cm]
		
		\textbf{Sinh Viên:}\\[0.3cm]
		Ngô Đức Phan Tiến Đạt\\[1cm]
		
		\textbf{Giảng viên hướng dẫn:}\\[0.3cm]
		Đoàn Thị Hồng Phước\\[1cm]
		
		\textbf{Năm:} 2026
		\vfill
	\end{titlepage}
	
	% ============================================================
	% LỜI CAM ĐOAN
	% ============================================================
	\thispagestyle{empty}
	\vspace*{3cm}
	\section*{LỜI CAM ĐOAN}
	
	Em xin cam đoan rằng khoá luận tốt nghiệp này là công trình nghiên cứu của
	bản thân, được thực hiện dưới sự hướng dẫn khoa học của giảng viên hướng dẫn
	tại Khoa Công nghệ Thông tin, Trường Đại Học Khoa Học -- Đại Học Huế. Toàn
	bộ nội dung trình bày là trung thực; mọi tham chiếu đến kết quả của các tác
	giả khác đều được ghi rõ nguồn gốc.
	
	\vspace{2cm}
	\begin{flushright}
		Thành phố Huế, ngày \quad tháng \quad năm 2026\\[1cm]
		\textbf{Ngô Đức Phan Tiến Đạt}
	\end{flushright}
	\clearpage
	
	% ============================================================
	% LỜI CẢM ƠN
	% ============================================================
	\thispagestyle{empty}
	\vspace*{3cm}
	\section*{LỜI CẢM ƠN}
	
	Em xin trân trọng cảm ơn cô Đoàn Thị Hồng Phước đã tận tình định hướng
	nghiên cứu và phản biện kỹ thuật trong suốt quá trình thực hiện khoá luận.
	Em cũng xin cảm ơn Khoa Công nghệ Thông tin, Trường Đại Học Khoa Học --
	Đại Học Huế đã tạo điều kiện để triển khai các thực nghiệm.
	
	\vspace{2cm}
	\begin{flushright}
		Thành phố Huế, ngày \quad tháng \quad năm 2026
	\end{flushright}
	\clearpage
	
	% ============================================================
	% TÓM TẮT
	% ============================================================
	\thispagestyle{empty}
	\vspace*{2cm}
	\section*{TÓM TẮT}
	\addcontentsline{toc}{section}{Tóm tắt}
	
	Khoá luận nghiên cứu bài toán phân đoạn văn bản (text chunking) trong kiến
	trúc Retrieval-Augmented Generation (RAG), với ứng dụng vào hệ thống tư vấn
	tuyển sinh Trường Đại Học Khoa Học -- Đại Học Huế (HUSC).
	
	\textbf{Câu hỏi nghiên cứu trung tâm:} Chiến lược chunking nào --- về mặt
	lý thuyết và thực nghiệm --- phù hợp nhất với đặc trưng văn bản hành chính
	tiếng Việt trong bài toán hỏi đáp tuyển sinh đại học?
	
	Khoá luận đóng góp trên ba phương diện. Thứ nhất, hệ thống hoá phân tích
	toán học bốn thế hệ chiến lược chunking: fixed-size (G1), semantic (G2),
	proposition/late chunking (G3), và RAPTOR/Contextual Retrieval (G4). Thứ
	hai, thiết lập khung đánh giá ba chỉ số (Recall@$k$, Precision@$k$, IoU)
	gắn với bộ dữ liệu tuyển sinh HUSC gồm 50 câu hỏi phân tầng. Thứ ba, đề
	xuất chiến lược lai (hybrid chunking) kết hợp semantic chunking cho tầng
	vector search với proposition extraction cho tầng xây dựng đồ thị tri thức.
	
	\textbf{Từ khóa:} Text chunking, Retrieval-Augmented Generation, Semantic
	Chunking, Late Chunking, RAPTOR, Proposition-based Retrieval, Tư vấn tuyển sinh.
	
	\clearpage
	
	% ============================================================
	% MỤC LỤC
	% ============================================================
	\tableofcontents
	\clearpage
	\listoffigures
	\addcontentsline{toc}{section}{Danh mục hình vẽ}
	\clearpage
	\listoftables
	\addcontentsline{toc}{section}{Danh mục bảng biểu}
	\clearpage
	
	% ============================================================
	% DANH MỤC KÝ HIỆU
	% ============================================================
	\section*{DANH MỤC KÝ HIỆU VÀ CHỮ VIẾT TẮT}
	\addcontentsline{toc}{section}{Danh mục ký hiệu và chữ viết tắt}
	
	\begin{table}[h]
		\centering
		\begin{tabular}{ll}
			\toprule
			\textbf{Ký hiệu / Viết tắt} & \textbf{Giải thích} \\
			\midrule
			RAG     & Retrieval-Augmented Generation \\
			LLM     & Large Language Model \\
			G1--G4  & Thế hệ chiến lược chunking 1--4 \\
			$D$     & Tài liệu đầu vào (chuỗi token) \\
			$C_i$   & Chunk thứ $i$ sau phân đoạn \\
			$\ell$  & Kích thước chunk (token) \\
			$s$     & Độ chồng lấp (overlap stride) \\
			$\phi$  & Hàm nhúng (embedding function) \\
			$S_C$   & Cosine similarity \\
			$\delta_{\text{thresh}}$ & Ngưỡng bất tương đồng ngữ nghĩa \\
			BM25    & Best Matching 25 (sparse retrieval) \\
			RRF     & Reciprocal Rank Fusion \\
			GMM     & Gaussian Mixture Model \\
			UMAP    & Uniform Manifold Approximation and Projection \\
			NER     & Named Entity Recognition \\
			HUSC    & Trường Đại Học Khoa Học -- Đại Học Huế \\
			\bottomrule
		\end{tabular}
	\end{table}
	\clearpage
	
	% ============================================================
	% PHẦN NỘI DUNG CHÍNH
	% ============================================================
	\pagenumbering{arabic}
	
	% ============================================================
	% CHƯƠNG 1: TỔNG QUAN RAG VÀ ĐỘNG LỰC NGHIÊN CỨU
	% ============================================================
	
	\chapter{Tổng quan về Retrieval-Augmented Generation và Vai trò của Chunking}
	\label{chap:rag_overview}
	
	\section{Hạn chế của Mô hình Ngôn ngữ Lớn trong Bài toán Hỏi đáp Chuyên biệt}
	\label{sec:llm_limits}
	
	Mô hình ngôn ngữ lớn (Large Language Model -- LLM) đã chứng minh khả năng
	xử lý ngôn ngữ tự nhiên vượt trội trên nhiều tác vụ tổng quát. Tuy nhiên,
	trong bối cảnh hệ thống hỏi đáp yêu cầu tri thức chuyên biệt, cập nhật và
	có thể kiểm chứng, LLM bộc lộ bốn hạn chế cấu trúc không thể khắc phục bằng
	cách tăng quy mô tham số~\cite{lewis2020retrieval}.
	
	\textbf{Ngưỡng kiến thức tĩnh (Static Knowledge Cutoff).} Quá trình tiền
	huấn luyện của LLM sử dụng tập dữ liệu thu thập đến một thời điểm cố định.
	Mọi thay đổi xảy ra sau đó --- chỉ tiêu tuyển sinh, mức học phí, quy chế
	xét tuyển --- đều nằm ngoài tầm nhận thức của mô hình, bất kể độ lớn kiến
	trúc.
	
	\textbf{Tri thức tổ chức nội bộ.} Văn bản tuyển sinh, quy chế nội bộ, và
	cơ sở dữ liệu sinh viên của HUSC chưa bao giờ xuất hiện trong bất kỳ tập
	huấn luyện công khai nào. LLM không thể lý luận trực tiếp trên những nguồn
	này.
	
	\textbf{Ảo giác (Hallucination).} LLM biểu diễn tri thức dưới dạng phân
	phối xác suất trên không gian ngôn ngữ, không phải tra cứu bản ghi. Do đó,
	mô hình có thể tự tin sản sinh thông tin sai lệch --- từ tên ngành học đến
	mã ngành và điểm chuẩn --- mà không phát tín hiệu cảnh báo.
	
	\textbf{Không thể kiểm chứng nguồn gốc.} Câu trả lời của LLM không gắn
	với bất kỳ đoạn văn bản cụ thể nào, khiến quá trình xác minh độ tin cậy
	trở nên không khả thi trong bối cảnh học thuật và hành chính.
	
	Phương án tinh chỉnh (fine-tuning) có vẻ trực quan nhưng không thực tế: chu
	kỳ tái huấn luyện tốn kém, không đảm bảo mô hình ghi nhớ thông tin theo cách
	có thể kiểm soát, và phải lặp lại mỗi khi tri thức thay đổi. RAG ra đời như
	một giải pháp bổ sung tri thức tại thời điểm suy luận (inference-time),
	tách biệt vòng đời của mô hình khỏi vòng đời của dữ liệu.
	
	% -------------------------------------------------------------------
	\section{Kiến trúc RAG: Pipeline Hình thức}
	\label{sec:rag_pipeline}
	
	Kiến trúc RAG tổng quát bao gồm ba thành phần tách biệt về mặt chức
	năng~\cite{lewis2020retrieval}. Gọi $\mathcal{D} = \{D_1, D_2, \ldots,
	D_M\}$ là tập tài liệu gốc và $q$ là truy vấn người dùng.
	
	\textbf{Thành phần 1 --- Indexing.} Mỗi tài liệu $D_j$ được xử lý thành
	tập chunk $\{C_1^j, C_2^j, \ldots, C_{m_j}^j\}$, sau đó mỗi chunk được
	nhúng vào không gian vector $\mathbb{R}^d$ qua hàm $\phi$:
	\begin{equation}
		\mathbf{v}_{C_i^j} = \phi(C_i^j) \in \mathbb{R}^d
		\label{eq:indexing}
	\end{equation}
	Toàn bộ $\sum_j m_j$ vector được lưu trong cơ sở dữ liệu vector (vector
	store) $\mathcal{V}$.
	
	\textbf{Thành phần 2 --- Retrieval.} Với truy vấn $q$, tìm $k$ chunk có
	độ tương đồng Cosine cao nhất:
	\begin{equation}
		S_C(\mathbf{v}_q, \mathbf{v}_{C}) =
		\frac{\mathbf{v}_q \cdot \mathbf{v}_{C}}
		{\|\mathbf{v}_q\|_2 \cdot \|\mathbf{v}_{C}\|_2}
		\label{eq:cosine}
	\end{equation}
	\begin{equation}
		\mathcal{R}(q, k) = \operatorname{Top-}k_{\,C \in \mathcal{V}}
		S_C\!\left(\phi(q),\, \phi(C)\right)
		\label{eq:retrieval}
	\end{equation}
	
	\textbf{Thành phần 3 --- Generation.} LLM sinh câu trả lời $\hat{a}$ dựa
	trên truy vấn và ngữ cảnh truy xuất được:
	\begin{equation}
		\hat{a} = \text{LLM}\!\left(q,\, \mathcal{R}(q, k)\right)
		\label{eq:generation}
	\end{equation}
	
	Điểm quan trọng cần lưu ý: bước Indexing --- đặc biệt là quyết định chia
	$D_j$ thành các $C_i^j$ như thế nào --- được thực hiện \emph{một lần, ngoại
		tuyến}, trong khi bước Retrieval và Generation được thực hiện \emph{trực
		tuyến} theo từng truy vấn. Do đó, chất lượng của chiến lược chunking ảnh
	hưởng đến mọi truy vấn trong suốt vòng đời của hệ thống.
	
	Hình~\ref{fig:rag_pipeline_full} minh họa toàn bộ pipeline.
	
	\begin{figure}[htbp]
		\centering
		\begin{adjustbox}{max width=\linewidth}
			\begin{tikzpicture}[
				box/.style={rectangle, rounded corners=4pt, draw=black!65,
					minimum height=0.85cm, align=center, font=\small,
					text width=2.4cm},
				arrow/.style={-Stealth, thick},
				phase/.style={rectangle, rounded corners=6pt, draw,
					minimum height=1.6cm, align=center, font=\footnotesize,
					text width=2.8cm, dashed}
				]
				% ---- Offline Phase ----
				\node[box, fill=blue!10] (docs)   at (0,   0) {Tài liệu\\gốc $\mathcal{D}$};
				\node[box, fill=blue!15] (chunk)  at (3.2, 0) {Chunking\\$\{C_i\}$};
				\node[box, fill=blue!20] (embed)  at (6.4, 0) {Embedding\\$\phi(C_i)$};
				\node[box, fill=blue!25] (store)  at (9.6, 0) {Vector\\Store $\mathcal{V}$};
				
				\draw[arrow] (docs)  -- node[above, font=\tiny]{Phân đoạn} (chunk);
				\draw[arrow] (chunk) -- node[above, font=\tiny]{Nhúng}     (embed);
				\draw[arrow] (embed) -- node[above, font=\tiny]{Lưu}       (store);
				
				\begin{scope}[on background layer]
					\node[draw=blue!40, fill=blue!4, rounded corners=8pt,
					fit=(docs)(chunk)(embed)(store),
					label={[font=\footnotesize\bfseries, text=blue!60]
						above:Giai đoạn Offline (Indexing)}] {};
				\end{scope}
				
				% ---- Online Phase ----
				\node[box, fill=orange!10] (query)  at (0,  -2.5) {Truy vấn\\người dùng $q$};
				\node[box, fill=orange!15] (ret)    at (3.2,-2.5) {Retrieval\\$\mathcal{R}(q,k)$};
				\node[box, fill=orange!20] (llm)    at (6.4,-2.5) {LLM\\Generation};
				\node[box, fill=green!20,
				draw=green!60]       (ans)    at (9.6,-2.5) {Câu trả\\lời $\hat{a}$};
				
				\draw[arrow] (query) -- (ret);
				\draw[arrow] (ret)   -- node[above, font=\tiny]{Ngữ cảnh} (llm);
				\draw[arrow] (llm)   -- (ans);
				
				% Kết nối vector store xuống retrieval
				\draw[arrow, gray!60, dashed] (store.south) -- ++(0,-0.4)
				-| (ret.north);
				
				\begin{scope}[on background layer]
					\node[draw=orange!40, fill=orange!4, rounded corners=8pt,
					fit=(query)(ret)(llm)(ans),
					label={[font=\footnotesize\bfseries, text=orange!70]
						below:Giai đoạn Online (per Query)}] {};
				\end{scope}
			\end{tikzpicture}
		\end{adjustbox}
		\caption{Pipeline RAG đầy đủ gồm hai giai đoạn. Giai đoạn Offline
			(trên) thực hiện một lần: phân đoạn, nhúng và lưu trữ. Giai đoạn
			Online (dưới) lặp lại theo từng truy vấn: truy xuất và tạo sinh.
			Chất lượng chunking ở giai đoạn Offline ảnh hưởng trực tiếp đến
			mọi truy vấn ở giai đoạn Online.}
		\label{fig:rag_pipeline_full}
	\end{figure}
	
	% -------------------------------------------------------------------
	\section{Chunking là Nút Thắt Cổ Chai của RAG}
	\label{sec:chunking_bottleneck}
	
	Barnett và cộng sự~\cite{barnett2024seven} tiến hành phân tích hệ thống bảy
	điểm thất bại phổ biến nhất khi triển khai RAG trong thực tế. Trong số đó,
	hai điểm liên quan trực tiếp đến chiến lược chunking và cùng nhau chiếm tỷ
	lệ thất bại cao nhất trong khảo sát thực tế:
	
	Điểm thất bại \#1 --- \emph{Missed Top Ranked Documents}: chunk liên quan
	bị thiếu trong tập top-$k$ vì ranh giới phân đoạn cắt đứt thông tin ngữ
	nghĩa. Điểm thất bại \#2 --- \emph{Not in Context}: chunk có nằm trong
	tập top-$k$ nhưng không được LLM sử dụng vì kích thước context window bị
	chiếm dụng bởi các chunk không liên quan truy xuất cùng lúc.
	
	Cả hai điểm thất bại này đều có căn nguyên toán học từ sự mất cân bằng
	giữa hai mục tiêu mâu thuẫn trong phân đoạn. Gọi $f(\ell)$ là Precision
	của truy xuất (chất lượng embedding cho chunk kích thước $\ell$) và
	$g(\ell)$ là Context Completeness (khả năng một chunk chứa đủ thông tin
	để trả lời câu hỏi). Hai hàm này thỏa mãn:
	\begin{equation}
		f'(\ell) < 0 \quad \text{và} \quad g'(\ell) > 0
		\qquad \forall\, \ell > 0
		\label{eq:chunking_tradeoff}
	\end{equation}
	Mối quan hệ đơn điệu nghịch chiều này --- đã được xác nhận thực nghiệm
	bởi Chroma Research~\cite{chroma2024eval} và Merola \& Singh~\cite{merola2025reconstructing}
	--- đặt ra bài toán tối ưu hóa đa mục tiêu không tầm thường. Hàm mục tiêu
	hợp nhất với trọng số $\lambda \in (0,1)$ phản ánh ưu tiên ứng dụng:
	\begin{equation}
		\mathcal{L}(\ell) = \lambda\, f(\ell) + (1-\lambda)\, g(\ell)
		\label{eq:chunking_objective}
	\end{equation}
	Điểm cực trị $\ell^*$ thỏa mãn điều kiện bậc nhất
	$\lambda f'(\ell^*) + (1-\lambda)g'(\ell^*) = 0$ và phụ thuộc vào phân
	phối nội dung corpus --- không tồn tại giá trị $\ell^*$ vạn năng áp dụng
	được cho mọi bài toán.
	
	% -------------------------------------------------------------------
	\section{Bài toán Tư vấn Tuyển sinh HUSC}
	\label{sec:husc_problem}
	
	Tư vấn tuyển sinh đại học là bài toán hỏi đáp có đặc trưng riêng biệt đặt
	ra những yêu cầu cụ thể cho thiết kế chunking.
	
	\textbf{Đặc trưng ngôn ngữ.} Văn bản tuyển sinh tiếng Việt tại HUSC có
	mật độ từ viết tắt cao (CNTT, ĐHKH, ĐH Huế), sử dụng tên riêng tổ chức
	có nhiều biến thể, và chứa các mệnh đề điều kiện phức hợp (``Thí sinh đạt
	từ 24,5 điểm trở lên theo tổ hợp A00, B00 hoặc D01 \emph{và} có điểm
	tiếng Anh riêng lẻ từ 7,0 điểm (thang 10) trở lên \ldots'').
	
	\textbf{Đặc trưng truy vấn.} Câu hỏi tuyển sinh phân bổ theo ba loại rõ
	nét, mỗi loại đặt ra yêu cầu khác nhau cho chiến lược chunking. Câu hỏi
	\emph{thực thể} (``Điểm chuẩn ngành CNTT năm 2024 là bao nhiêu?'') yêu
	cầu truy xuất độ chính xác cao, hưởng lợi từ chunk nhỏ. Câu hỏi
	\emph{đa điều kiện} (``Tôi có $23$ điểm thi tốt nghiệp, học lực giỏi, và
	có chứng chỉ IELTS~6.5, có được xét tuyển không?'') đòi hỏi tổng hợp
	thông tin từ nhiều đoạn văn bản. Câu hỏi \emph{so sánh} (``Chương trình
	tiêu chuẩn và chương trình chất lượng cao của ngành CNTT khác nhau như
	thế nào?'') yêu cầu truy xuất ở mức độ trừu tượng cao.
	
	\textbf{Quy mô dữ liệu.} Corpus tuyển sinh HUSC bao gồm các tài liệu:
	Đề án tuyển sinh hàng năm (khoảng 80--120 trang), Quy chế tuyển sinh Bộ
	GD\&ĐT (khoảng 40--60 trang), trang web tuyển sinh chính thức, và các
	thông báo bổ sung. Quy mô trung bình: $N \approx 120{,}000$ token toàn
	corpus, với cấu trúc hành chính gồm chương mục, bảng số liệu, và mệnh đề
	điều kiện.
	
	\begin{table}[htbp]
		\centering
		\caption{Phân tích đặc trưng câu hỏi tuyển sinh HUSC theo loại truy vấn}
		\label{tab:husc_query_types}
		\begin{tabular}{lccl}
			\toprule
			\textbf{Loại truy vấn} & \textbf{Tỷ lệ (\%)} & \textbf{Độ phức tạp}
			& \textbf{Yêu cầu chunking} \\
			\midrule
			Thực thể đơn        & 40 & Thấp    & Chunk nhỏ, nguyên tử \\
			Đa điều kiện        & 35 & Trung bình & Chunk trung bình, đủ ngữ cảnh \\
			So sánh / Tổng hợp  & 25 & Cao     & Đa cấp độ trừu tượng \\
			\bottomrule
		\end{tabular}
	\end{table}
	
	\noindent\textit{Lưu ý: Tỷ lệ trong Bảng~\ref{tab:husc_query_types} là ước
		lượng dựa trên phân tích định tính bộ câu hỏi thực nghiệm; chưa có bằng
		chứng định lượng từ nghiên cứu đã xuất bản về phân phối câu hỏi tuyển sinh
		tiếng Việt.}
	
	Sự đa dạng này là lý do chính mà không có chiến lược chunking đơn lẻ nào
	tối ưu đồng thời cho cả ba loại truy vấn. Nghiên cứu trong khoá luận này
	nhằm xác định chiến lược (hoặc tổ hợp chiến lược) phù hợp nhất với phân
	phối truy vấn cụ thể của bài toán HUSC.
	
	% ============================================================
	% CHƯƠNG 2: LÝ THUYẾT CHUNKING (G1–G4)
	% ============================================================
	
	\chapter{Lý thuyết Phân đoạn Văn bản: Từ G1 đến G4}
	\label{chap:chunking_theory}
	
	\section{Định nghĩa Hình thức Bài toán Chunking}
	\label{sec:chunking_formal}
	
	\begin{dinhnghia}[Phân đoạn văn bản]
		\label{def:chunking}
		Cho tài liệu $D = (t_1, t_2, \ldots, t_N)$ gồm $N$ token và tham số
		$\ell_{\max} > 0$. Một \emph{phân đoạn} (chunking) của $D$ là ánh xạ
		\[
		\pi(D) = (C_1, C_2, \ldots, C_m)
		\]
		trong đó:
		\begin{enumerate}[label=(\roman*)]
			\item $C_i = (t_{a_i}, t_{a_i+1}, \ldots, t_{b_i})$ là chuỗi token
			liên tiếp với $1 \leq a_i \leq b_i \leq N$;
			\item $|C_i| = b_i - a_i + 1 \leq \ell_{\max}$ với mọi $i$;
			\item $\bigcup_{i=1}^{m} \{t_{a_i}, \ldots, t_{b_i}\} = \{t_1, \ldots, t_N\}$
			(tính phủ toàn bộ).
		\end{enumerate}
		Khi cho phép chồng lấp kiểm soát, thay điều kiện (i) bằng yêu cầu mỗi
		token xuất hiện trong ít nhất một chunk.
	\end{dinhnghia}
	
	Hàm chất lượng của một phân đoạn $\pi$ trên bộ câu hỏi $\mathcal{Q}$:
	\begin{equation}
		Q(\pi, \mathcal{Q}) = \frac{1}{|\mathcal{Q}|}
		\sum_{q \in \mathcal{Q}} \text{Score}\!\left(
		\mathcal{R}(q, k; \pi),\; a^*_q
		\right)
		\label{eq:chunking_quality}
	\end{equation}
	trong đó $\mathcal{R}(q, k; \pi)$ là tập $k$ chunk được truy xuất khi áp
	dụng phân đoạn $\pi$, và $a^*_q$ là câu trả lời chuẩn. Hàm $\text{Score}$
	cụ thể sẽ được định nghĩa tại Chương~\ref{chap:experiments}.
	
	% -------------------------------------------------------------------
	\section{Thế hệ G1: Phân đoạn theo Cấu trúc Hình thức}
	\label{sec:g1}
	
	\subsection{Fixed-size Chunking}
	\label{subsec:fixed_size}
	
	Fixed-size chunking là chiến lược đơn giản nhất và là baseline trong phần
	lớn tài liệu kỹ thuật~\cite{lewis2020retrieval}. Tài liệu được cắt thành
	các đoạn có độ dài $\ell$ token với stride $s < \ell$ (cho phép chồng lấp):
	\begin{equation}
		C_i = (t_{(i-1)(\ell-s)+1},\; t_{(i-1)(\ell-s)+2},\; \ldots,\;
		t_{(i-1)(\ell-s)+\ell})
		\label{eq:fixed_chunks}
	\end{equation}
	Số chunk tạo ra:
	\begin{equation}
		m = \left\lceil \frac{N - s}{\ell - s} \right\rceil
		\label{eq:chunk_count}
	\end{equation}
	Trường hợp đặc biệt $s = 0$ (không chồng lấp) cho $m = \lceil N/\ell \rceil$.
	
	\textbf{Hiệu ứng biên (Boundary Effect).} Xác suất một ranh giới cắt ngẫu
	nhiên nằm giữa câu hoàn chỉnh tỷ lệ nghịch với độ dài câu trung bình
	$\bar{\ell}_s$ của corpus:
	\begin{equation}
		P(\text{cắt giữa câu}) \approx 1 - \frac{\bar{\ell}_s}{\ell}
		\quad \text{với } \ell \gg \bar{\ell}_s
		\label{eq:boundary_prob}
	\end{equation}
	Với corpus văn bản tuyển sinh tiếng Việt ($\bar{\ell}_s \approx 18$--$25$
	token) và $\ell = 200$ token, xác suất này xấp xỉ $90$--$91\%$ --- nghĩa
	là 9 trong 10 chunk có khả năng chứa câu bị cắt đứt ở một trong hai đầu.
	
	\subsection{Recursive Chunking}
	\label{subsec:recursive}
	
	Recursive Character Text Splitter (phổ biến trong LangChain) giải quyết
	hiệu ứng biên bằng cách định nghĩa danh sách ưu tiên tách
	$\Sigma = [\sigma_1, \sigma_2, \ldots, \sigma_K]$ và áp dụng đệ quy:
	\begin{equation}
		\text{Split}(D, \Sigma, \ell) =
		\begin{cases}
			\{D\} & \text{nếu } |D| \leq \ell \\[4pt]
			\displaystyle\bigcup_{c \,\in\, T(\sigma_1, D)}
			\text{Split}(c, \Sigma, \ell)
			& \text{nếu } |D| > \ell
		\end{cases}
		\label{eq:recursive_split}
	\end{equation}
	trong đó $T(\sigma, D)$ là hàm tách $D$ tại ký tự $\sigma$. Nếu tách theo
	$\sigma_1$ không tạo ra chunk nhỏ hơn $\ell$, thuật toán thử $\sigma_2$,
	$\sigma_3$, \ldots theo thứ tự. Thứ tự ưu tiên điển hình:
	$\Sigma = [\texttt{\textquotedbl\textbackslash n\textbackslash n\textquotedbl},\;
	\texttt{\textquotedbl\textbackslash n\textquotedbl},\;
	\texttt{\textquotedbl.\textquotedbl},\;
	\texttt{\textquotedbl,\textquotedbl},\;
	\texttt{\textquotedbl\ \textquotedbl}]$.
	
	Recursive chunking vô hiệu hóa hiệu ứng biên ở mức đoạn văn và câu trong
	phần lớn trường hợp. Tuy nhiên, hai chiến lược G1 chia sẻ hạn chế cấu trúc
	cơ bản: cả hai xử lý văn bản như \emph{chuỗi ký tự} mà không có bất kỳ
	biểu diễn ngữ nghĩa nào. Ranh giới phân đoạn không phản ánh cấu trúc ý
	nghĩa của tài liệu.
	
	% -------------------------------------------------------------------
	\section{Thế hệ G2: Phân đoạn theo Ngữ nghĩa Embedding}
	\label{sec:g2}
	
	\subsection{Phát hiện Ranh giới Ngữ nghĩa}
	\label{subsec:semantic_boundary}
	
	Semantic chunking (Kamradt, 2024) bổ sung một lớp nhận thức ngữ nghĩa:
	thay vì dựa vào ký tự phân tách, chiến lược này phát hiện điểm đứt gãy
	ngữ nghĩa (semantic breakpoint) dựa trên độ tương đồng embedding giữa các
	câu liền kề.
	
	Gọi $\{s_1, s_2, \ldots, s_n\}$ là chuỗi câu trong tài liệu và
	$\mathbf{v}_i = \phi(s_i)$ là embedding của câu $s_i$. Định nghĩa khoảng
	cách ngữ nghĩa tại vị trí $i$:
	\begin{equation}
		d_i = 1 - S_C(\mathbf{v}_i,\, \mathbf{v}_{i+1}), \quad i = 1, \ldots, n-1
		\label{eq:semantic_dist}
	\end{equation}
	Ranh giới chunk được đặt tại vị trí $i^*$ khi:
	\begin{equation}
		d_{i^*} \geq \delta_{\text{thresh}}
		\label{eq:semantic_breakpoint}
	\end{equation}
	Ngưỡng $\delta_{\text{thresh}}$ được xác định tự thích nghi theo phân vị
	của phân phối khoảng cách:
	\begin{equation}
		\delta_{\text{thresh}} = F^{-1}_{d}(p)
		\quad \text{với } p \in [0.90, 0.97]
		\label{eq:threshold_adaptive}
	\end{equation}
	trong đó $F^{-1}_d$ là hàm phân vị của $\{d_i\}$. Thực tiễn, $p = 0.95$
	(phân vị thứ 95) cho kết quả ổn định trên nhiều corpus~\cite{chroma2024eval}.
	
	\subsection{Parent Document Retriever: Tách Biệt Đơn vị Truy xuất và Ngữ cảnh}
	\label{subsec:parent_doc}
	
	Parent Document Retriever giải quyết mâu thuẫn ở~\eqref{eq:chunking_tradeoff}
	bằng cách duy trì \emph{hai cấp biểu diễn song song}. Child chunks
	$C_i^{\text{child}}$ có kích thước nhỏ ($\ell_c \approx 100$--$200$ token)
	được dùng để lập chỉ mục và truy xuất chính xác. Khi chunk con được xác
	định, hệ thống trả về parent document $C_j^{\text{par}}$ --- đoạn cha lớn
	hơn bao hàm chunk con đó:
	\begin{equation}
		\text{Context}(q) = C_j^{\text{par}}, \quad
		j = \operatorname{parent}(i^*)
		\label{eq:parent_doc}
	\end{equation}
	trong đó $i^* = \argmax_{i} S_C(\phi(q), \phi(C_i^{\text{child}}))$.
	Kỹ thuật này thừa nhận một nguyên lý quan trọng: đơn vị tối ưu cho
	\emph{truy xuất} và đơn vị tối ưu cho \emph{tạo sinh} không nhất thiết
	trùng nhau.
	
	\subsection{Hạn chế Cấu trúc của G1 và G2}
	\label{subsec:early_chunking_limit}
	
	Cả G1 và G2 đều thuộc kiến trúc \emph{early chunking}: tài liệu được phân
	đoạn \emph{trước khi} đưa vào encoder. Embedding của chunk $C_i$ là:
	\begin{equation}
		\mathbf{v}_{C_i}^{\text{early}} = \text{Pool}\!\left(E(C_i)\right)
		\label{eq:early_embed}
	\end{equation}
	trong đó $E(C_i)$ chỉ nhận $C_i$ làm đầu vào. Do cơ chế Self-Attention
	của Transformer~\cite{vaswani2017attention}, mỗi token trong $C_i$ chỉ
	attend đến các token khác trong $C_i$, không thấy ngữ cảnh bên ngoài chunk.
	
	Hậu quả là \emph{mất ngữ cảnh nội chiếu} (coreference loss): các đại từ,
	mệnh đề tham chiếu và cụm từ tỉnh lược bên trong $C_i$ sẽ bị biểu diễn
	sai. Günther và cộng sự~\cite{gunther2024late} minh họa bằng ví dụ cụ thể:
	trong tài liệu tuyển sinh, câu ``\textit{Ngưỡng điểm đảm bảo chất lượng
		đầu vào là 18,0 điểm.}'' được phân đoạn tách khỏi đoạn mở đầu xác định
	rõ đây là quy định của ngành Công nghệ Thông tin. Với early chunking,
	embedding của chunk câu đó có cosine thấp với truy vấn ``điểm chuẩn CNTT''
	vì nhãn định danh ngành không xuất hiện trong chunk đó.
	
	% -------------------------------------------------------------------
	\section{Thế hệ G3: Phân đoạn sau Encoder}
	\label{sec:g3}
	
	\subsection{Late Chunking: Đảo ngược Thứ tự Xử lý}
	\label{subsec:late_chunking}
	
	Late Chunking~\cite{gunther2024late} (Günther et al., 2024, Jina AI) đảo
	ngược thứ tự truyền thống: thay vì \emph{phân đoạn trước, mã hóa sau},
	late chunking \emph{mã hóa toàn tài liệu trước, phân đoạn sau}.
	
	Gọi $E$ là bộ mã hóa long-context với cửa sổ ngữ cảnh $W$. Giả sử
	$N \leq W$ (tài liệu vừa trong cửa sổ ngữ cảnh). Embedding của chunk
	$C_i = (t_{a_i}, \ldots, t_{b_i})$ trong kiến trúc late chunking là:
	\begin{equation}
		\mathbf{v}_{C_i}^{\text{late}}
		= \frac{1}{b_i - a_i + 1} \sum_{j=a_i}^{b_i} E_j(D)
		\label{eq:late_embed}
	\end{equation}
	trong đó $E_j(D)$ là contextualized token embedding của token $t_j$ khi
	\emph{toàn bộ tài liệu $D$} được đưa vào encoder. Khác với~\eqref{eq:early_embed}
	trong đó $E_j(C_i)$ chỉ thấy $C_i$: do cơ chế Self-Attention, $E_j(D)$
	đã ``nhìn thấy'' mọi token trong $D$ thông qua attention weight
	$\alpha_{jk} = \text{softmax}(\mathbf{q}_j \cdot \mathbf{k}_k / \sqrt{d_h})$
	với mọi $k \in [1, N]$. Do đó $\mathbf{v}_{C_i}^{\text{late}}$ mang thông
	tin ngữ cảnh toàn cục ngay cả khi $C_i$ chứa đại từ hay tham chiếu ngoài biên.
	
	\begin{figure}[htbp]
		\centering
		\begin{adjustbox}{max width=\linewidth}
			\begin{tikzpicture}[
				box/.style={rectangle, draw=black!60, minimum height=0.65cm,
					align=center, font=\small, rounded corners=2pt},
				arrow/.style={-Stealth, thick}
				]
				% ---- EARLY ----
				\node[font=\small\bfseries] at (-3.5, 0.7) {Early Chunking};
				\node[box, fill=gray!12, text width=5cm] (D_e) at (-3.5, 0)
				{Tài liệu $D = (t_1 \ldots t_N)$};
				\node[box, fill=blue!10, text width=2cm] (C1_e) at (-5.2,-1.3)
				{$E(C_1)$\\{\tiny chỉ thấy $C_1$}};
				\node[box, fill=blue!10, text width=2cm] (C2_e) at (-1.8,-1.3)
				{$E(C_2)$\\{\tiny chỉ thấy $C_2$}};
				\node[box, fill=green!10, text width=2cm] (v1_e) at (-5.2,-2.6)
				{$\mathbf{v}_{C_1}^{\text{early}}$};
				\node[box, fill=green!10, text width=2cm] (v2_e) at (-1.8,-2.6)
				{$\mathbf{v}_{C_2}^{\text{early}}$};
				\draw[arrow] (D_e.south) -- ++(0,-0.2) -| (C1_e.north);
				\draw[arrow] (D_e.south) -- ++(0,-0.2) -| (C2_e.north);
				\draw[arrow] (C1_e) -- (v1_e);
				\draw[arrow] (C2_e) -- (v2_e);
				\node[font=\tiny, red!70] at (-3.5,-1.6)
				{\textit{Không thấy nhau}};
				
				% ---- LATE ----
				\node[font=\small\bfseries] at (4.5, 0.7) {Late Chunking};
				\node[box, fill=gray!12, text width=5.5cm] (D_l) at (4.5, 0)
				{Tài liệu $D$ (toàn bộ $\leq W$)};
				\node[box, fill=blue!20, text width=5cm] (E_l) at (4.5,-1.3)
				{$E(D)$: long-context encoder\\{\tiny toàn ngữ cảnh qua Self-Attention}};
				\node[box, fill=green!10, text width=2cm] (v1_l) at (2.7,-2.6)
				{Slice $[a_1..b_1]$\\$\mathbf{v}_{C_1}^{\text{late}}$};
				\node[box, fill=green!10, text width=2cm] (v2_l) at (6.3,-2.6)
				{Slice $[a_2..b_2]$\\$\mathbf{v}_{C_2}^{\text{late}}$};
				\draw[arrow] (D_l) -- (E_l);
				\draw[arrow] (E_l.south) -- ++(0,-0.2) -| (v1_l.north);
				\draw[arrow] (E_l.south) -- ++(0,-0.2) -| (v2_l.north);
			\end{tikzpicture}
		\end{adjustbox}
		\caption{So sánh kiến trúc early chunking (trái) và late chunking (phải).
			Early chunking: mỗi chunk được mã hóa độc lập, token chỉ attend
			trong biên chunk. Late chunking: toàn tài liệu được mã hóa một lần
			qua long-context encoder; embedding của từng chunk được cắt ra từ
			chuỗi token embedding đã mang thông tin toàn cục.}
		\label{fig:late_vs_early}
	\end{figure}
	
	\textbf{Ràng buộc và phức tạp tính toán.} Ràng buộc $N \leq W$ yêu cầu bộ
	mã hóa long-context; các mô hình hiện đại như \texttt{jina-embeddings-v3}
	hỗ trợ $W = 8{,}192$ token, đủ cho hầu hết tài liệu tuyển sinh đơn lẻ.
	Độ phức tạp tính toán tăng từ $\mathcal{O}(m \cdot \ell^2)$ (early, xử lý
	$m$ chunk riêng lẻ) lên $\mathcal{O}(N^2)$ (late, một lần mã hóa toàn
	tài liệu). Tuy nhiên do chỉ thực hiện một lần trong giai đoạn offline,
	chi phí tăng thêm này thường chấp nhận được trong thực tế~\cite{gunther2024late}.
	
	\subsection{Proposition-based Chunking: Đơn vị Nguyên tử}
	\label{subsec:proposition}
	
	Dense X Retrieval~\cite{chen2024dense} (Chen et al., EMNLP 2024, trang
	15159--15177) giới thiệu đơn vị truy xuất hoàn toàn mới: \emph{proposition}
	(mệnh đề nguyên tử).
	
	\begin{dinhnghia}[Proposition]
		\label{def:proposition}
		Một proposition là biểu thức ngôn ngữ tự nhiên thỏa mãn ba điều kiện đồng
		thời:
		\begin{enumerate}[label=(\roman*)]
			\item \textbf{Tính nguyên tử:} chứa đúng một factoid độc lập, không
			thể phân tách thêm mà không mất nghĩa;
			\item \textbf{Tính tự chứa:} có thể hiểu đúng mà không cần tham chiếu
			đến ngữ cảnh ngoài;
			\item \textbf{Tính súc tích:} loại bỏ mọi thông tin dư thừa không ảnh
			hưởng đến factoid cốt lõi.
		\end{enumerate}
	\end{dinhnghia}
	
	Quá trình sinh proposition từ đoạn văn $P$ được thực hiện qua mô hình
	Propositionizer --- Flan-T5-large tinh chỉnh trên 42,857 cặp ví dụ sinh
	bởi GPT-4~\cite{chen2024dense}:
	\begin{equation}
		(p_1, p_2, \ldots, p_r) = \text{Propositionizer}(P)
		\label{eq:propositionizer}
	\end{equation}
	Mỗi $p_j$ là một câu hoàn chỉnh, tự chứa, trung bình dài $10$--$30$ token.
	
	\textbf{Kết quả thực nghiệm đã công bố.} Chen và cộng sự tiến hành so sánh
	trên năm tập hỏi đáp mở (NQ, TriviaQA, WebQ, SQuAD, EntityQuestions) với
	ba đơn vị truy xuất: passage ($100$ từ), sentence, và proposition. Kết quả
	chính với unsupervised dense retrievers:
	\begin{equation}
		\Delta \text{Recall}@20(\text{prop vs.\ passage}) = +10.1\%
		\text{ (trung bình 5 tập)}
		\label{eq:prop_recall20}
	\end{equation}
	Với supervised retrievers (đã huấn luyện trên passage-query pairs):
	\begin{equation}
		\Delta \text{Recall}@20(\text{prop vs.\ passage}) = +2.7\%
		\text{ (trung bình 5 tập)}
		\label{eq:prop_recall20_sup}
	\end{equation}
	Lợi thế của proposition tăng mạnh với long-tail entities: trên tập
	EntityQuestions với unsupervised retrievers, mức cải thiện Recall@5 đạt
	$17$--$25\%$ theo tương đối~\cite{chen2024dense}. Giải thích toán học:
	khi tài liệu được phân thành propositions, mỗi vector nhúng tập trung biểu
	diễn một factoid duy nhất, giảm đáng kể nhiễu attention từ các ý tưởng
	thứ yếu trong cùng passage.
	
	\begin{vidu}[Proposition trong văn bản tuyển sinh]
		Đoạn gốc: ``Ngành Công Nghệ Thông Tin (CNTT) có chỉ tiêu 60 sinh viên cho
		chương trình đại trà và 20 sinh viên cho chương trình chất lượng cao. Điểm
		chuẩn năm 2024 là 24,5 theo tổ hợp A00.''
		
		Sau propositionization:
		\begin{enumerate}
			\item ``Ngành Công Nghệ Thông Tin (CNTT) có chỉ tiêu 60 sinh viên
			cho chương trình đại trà.''
			\item ``Ngành Công Nghệ Thông Tin (CNTT) có chỉ tiêu 20 sinh viên
			cho chương trình chất lượng cao.''
			\item ``Điểm chuẩn ngành Công Nghệ Thông Tin năm 2024 là 24,5 theo
			tổ hợp A00.''
		\end{enumerate}
		Mỗi proposition chứa tên ngành tường minh --- tránh hoàn toàn hiệu ứng
		coreference loss mô tả tại Mục~\ref{subsec:early_chunking_limit}.
	\end{vidu}
	
	% -------------------------------------------------------------------
	\section{Thế hệ G4: Phân đoạn Đa cấp và Tăng cường Ngữ cảnh}
	\label{sec:g4}
	
	\subsection{RAPTOR: Cây Tóm tắt Đệ quy}
	\label{subsec:raptor}
	
	RAPTOR~\cite{sarthi2024raptor} (Sarthi et al., ICLR 2024) tiếp cận bài toán
	chunking từ góc độ hoàn toàn khác: thay vì tìm ranh giới phân đoạn tốt hơn,
	RAPTOR xây dựng một cây tóm tắt đa cấp cho phép truy xuất đồng thời ở nhiều
	mức trừu tượng.
	
	\textbf{Xây dựng cây (từ dưới lên).} Bước 1: Leaf nodes là các chunk văn
	bản gốc; mỗi chunk được nhúng và giảm chiều qua UMAP xuống không gian thấp
	chiều. Bước 2: Phân cụm bằng Gaussian Mixture Model (GMM):
	\begin{equation}
		p(\mathbf{v}) = \sum_{k=1}^{K} \pi_k\,
		\mathcal{N}(\mathbf{v};\, \boldsymbol{\mu}_k,\, \boldsymbol{\Sigma}_k)
		\label{eq:gmm}
	\end{equation}
	Số cụm $K$ được xác định bằng BIC (Bayesian Information Criterion). Bước
	3: Mỗi cụm được tóm tắt thành node cha bằng LLM. Bước 4: Lặp lại bước 1--3
	trên các node cha cho đến khi chỉ còn một root node.
	
	\textbf{Chiến lược truy xuất Tree Collapse.} Tại thời điểm truy vấn, toàn
	bộ node từ mọi cấp $l = 0, 1, \ldots, L$ được đưa vào một flat index duy
	nhất:
	\begin{equation}
		\mathcal{R}_{\text{RAPTOR}}(q, k) =
		\operatorname{Top-}k_{\,n \in \bigcup_{l=0}^L T_l}
		S_C(\phi(q),\, \phi(n))
		\label{eq:raptor_collapse}
	\end{equation}
	Cơ chế này cho phép một truy vấn trả về đồng thời cả leaf nodes (chi tiết
	cụ thể) và summary nodes (bức tranh tổng quan), giải quyết bài toán câu
	hỏi so sánh phân loại tại Bảng~\ref{tab:husc_query_types}.
	
	\textbf{Kết quả thực nghiệm đã công bố.} Kết hợp RAPTOR retrieval với
	GPT-4 cải thiện kết quả tốt nhất trước đó trên QuALITY benchmark (đánh
	giá hiểu văn bản dài) thêm $20\%$ (tuyệt đối)~\cite{sarthi2024raptor}.
	
	\subsection{Contextual Retrieval: Bổ sung Ngữ cảnh Tường minh}
	\label{subsec:contextual_retrieval}
	
	Contextual Retrieval~\cite{anthropic2024contextual} (Anthropic, 2024) tiếp
	cận vấn đề từ hướng khác: bổ sung ngữ cảnh một cách tường minh cho mỗi
	chunk trước khi nhúng, thay vì thay đổi kiến trúc encoder.
	
	Với mỗi chunk $C_i$, LLM sinh đoạn mô tả ngữ cảnh $\kappa_i$ dài
	$50$--$100$ token, đặt chunk vào bối cảnh tài liệu gốc:
	\begin{equation}
		\kappa_i = \text{LLM}(D,\; C_i,\; \texttt{prompt}_{\text{ctx}})
		\label{eq:contextualize}
	\end{equation}
	Chunk được làm giàu: $C_i^+ = \kappa_i \oplus C_i$. Toàn bộ corpus sau
	đó được lập chỉ mục theo chiến lược lai BM25 + dense, kết quả hợp nhất
	bằng RRF:
	\begin{equation}
		\text{RRF}(d) = \frac{w_{\text{BM25}}}{k_0 + r_{\text{BM25}}(d)}
		+ \frac{w_{\text{dense}}}{k_0 + r_{\text{dense}}(d)}
		\label{eq:rrf}
	\end{equation}
	
	Chi phí của Contextual Retrieval tập trung hoàn toàn ở giai đoạn indexing
	offline: mỗi chunk yêu cầu một lời gọi LLM với toàn tài liệu làm ngữ cảnh.
	Merola \& Singh~\cite{merola2025reconstructing} ghi nhận mức sử dụng VRAM
	$\approx 20$~GB trên RTX~4090 khi xử lý tập NFCorpus, khẳng định đây là
	chiến lược G4 tốn kém tài nguyên nhất nhưng không ảnh hưởng đến latency
	truy vấn trực tuyến.
	
	% -------------------------------------------------------------------
	\section{Phân tích So sánh và Tổng kết Lý thuyết}
	\label{sec:theory_comparison}
	
	Bảng~\ref{tab:chunking_comparison_full} tổng hợp đặc trưng của bốn thế hệ
	chiến lược. Hình~\ref{fig:pareto_frontier} biểu diễn đường cận Pareto trong
	không gian thiết kế (chi phí -- chất lượng ngữ cảnh).
	
	\begin{table}[htbp]
		\centering
		\caption{So sánh bốn thế hệ chiến lược chunking. Ký hiệu: $\checkmark$ tốt,
			$\sim$ trung bình, $\times$ kém, (?) cần đánh giá thực nghiệm.}
		\label{tab:chunking_comparison_full}
		\begin{adjustbox}{max width=\linewidth}
			\begin{tabular}{lccccc}
				\toprule
				\textbf{Chiến lược}
				& \textbf{Ngữ cảnh toàn cục}
				& \textbf{Factoid}
				& \textbf{Câu hỏi tổng hợp}
				& \textbf{Chi phí Index}
				& \textbf{Yêu cầu} \\
				\midrule
				Fixed-size (G1)      & $\times$     & $\sim$     & $\times$     & Rất thấp & Không \\
				Recursive (G1)       & $\times$     & $\sim$     & $\times$     & Thấp     & Không \\
				Semantic (G2)        & $\sim$       & $\checkmark$ & $\sim$     & Trung bình & Embedding \\
				Parent Doc (G2)      & $\sim$       & $\checkmark$ & $\sim$     & Thấp     & Không \\
				Late Chunking (G3)   & $\checkmark$ & $\checkmark$ & $\sim$     & Trung bình & Long-ctx encoder \\
				Proposition (G3)     & $\checkmark$ & $\checkmark$ & $\sim$     & Cao      & LLM fine-tuned \\
				RAPTOR (G4)          & $\checkmark$ & $\sim$     & $\checkmark$ & Cao      & LLM + GMM \\
				Contextual (G4)      & $\checkmark$ & $\checkmark$ & $\checkmark$ & Rất cao & LLM + GPU \\
				\bottomrule
			\end{tabular}
		\end{adjustbox}
	\end{table}
	
	\begin{figure}[htbp]
		\centering
		\begin{tikzpicture}
			\begin{axis}[
				width=10cm, height=7cm,
				xlabel={Chi phí tính toán (tương đối)},
				ylabel={Chất lượng bảo toàn ngữ cảnh},
				xlabel style={font=\small},
				ylabel style={font=\small},
				xmin=0.5, xmax=5.2,
				ymin=0.5, ymax=5.0,
				xtick={1,2,3,4,5},
				xticklabels={Rất thấp, Thấp, Trung bình, Cao, Rất cao},
				ytick={1,2,3,4},
				yticklabels={Rất thấp, Thấp, Trung bình, Cao},
				xticklabel style={font=\footnotesize},
				yticklabel style={font=\footnotesize},
				grid=major,
				grid style={draw=gray!20},
				legend style={at={(0.03,0.97)}, anchor=north west, font=\footnotesize},
				]
				\addplot[only marks, mark=square*, mark size=4pt,
				mark options={fill=gray!60}]
				coordinates {(1.0,1.2)(1.6,1.7)};
				\addlegendentry{G1}
				
				\addplot[only marks, mark=triangle*, mark size=4pt,
				mark options={fill=blue!60}]
				coordinates {(2.3,2.4)(1.9,2.1)};
				\addlegendentry{G2}
				
				\addplot[only marks, mark=diamond*, mark size=4pt,
				mark options={fill=orange!70}]
				coordinates {(2.9,3.2)(3.4,3.5)};
				\addlegendentry{G3}
				
				\addplot[only marks, mark=*, mark size=4pt,
				mark options={fill=red!70}]
				coordinates {(4.1,4.2)(4.8,4.6)};
				\addlegendentry{G4}
				
				\node[font=\tiny, right] at (axis cs:1.0,1.2) {Fixed};
				\node[font=\tiny, right] at (axis cs:1.6,1.7) {Recursive};
				\node[font=\tiny, right] at (axis cs:2.3,2.4) {Semantic};
				\node[font=\tiny, right] at (axis cs:1.9,2.1) {Parent};
				\node[font=\tiny, right] at (axis cs:2.9,3.2) {Late};
				\node[font=\tiny, right] at (axis cs:3.4,3.5) {Proposition};
				\node[font=\tiny, left]  at (axis cs:4.1,4.2) {RAPTOR};
				\node[font=\tiny, left]  at (axis cs:4.8,4.6) {Contextual};
				
				\addplot[dashed, gray!60, thick, smooth]
				coordinates {(1.0,1.2)(1.6,1.7)(2.3,2.4)(2.9,3.2)(4.1,4.2)(4.8,4.6)};
			\end{axis}
		\end{tikzpicture}
		\caption{Đường cận Pareto (đứt nét) trong không gian thiết kế chunking.
			Không có chiến lược nào vượt trội đồng thời trên cả hai chiều. Vị trí
			các điểm là ước lượng định tính tổng hợp từ
			\cite{chroma2024eval, merola2025reconstructing}; không phản ánh số liệu
			thực nghiệm tuyệt đối trên một tập dữ liệu cụ thể nào.}
		\label{fig:pareto_frontier}
	\end{figure}
	
	% ============================================================
	% CHƯƠNG 3: THIẾT KẾ THỰC NGHIỆM VÀ KẾT QUẢ
	% ============================================================
	
	\chapter{Thiết kế Thực nghiệm và Đánh giá}
	\label{chap:experiments}
	
	\section{Khung Đánh giá}
	\label{sec:eval_framework}
	
	Để so sánh các chiến lược chunking một cách công bằng và có thể tái lập,
	khoá luận áp dụng khung đánh giá theo ba chỉ số độc lập, mỗi chỉ số phản
	ánh một tầng khác nhau trong pipeline RAG.
	
	\subsection{Chỉ số Tầng Truy xuất}
	\label{subsec:retrieval_metrics}
	
	\textbf{Recall@$k$} đo tỷ lệ câu trả lời chuẩn được tìm thấy trong tập
	$k$ chunk truy xuất hàng đầu. Gọi $\mathcal{G}_q$ là tập token gold cần
	thiết và $\mathcal{R}(q,k;\pi)$ là tập token trong $k$ chunk được truy xuất:
	\begin{equation}
		\text{Recall}@k(q;\pi) =
		\frac{|\mathcal{R}(q,k;\pi) \cap \mathcal{G}_q|}{|\mathcal{G}_q|}
		\label{eq:recall_k}
	\end{equation}
	
	\textbf{Precision@$k$}:
	\begin{equation}
		\text{Precision}@k(q;\pi) =
		\frac{|\mathcal{R}(q,k;\pi) \cap \mathcal{G}_q|}
		{|\mathcal{R}(q,k;\pi)|}
		\label{eq:precision_k}
	\end{equation}
	
	\textbf{Intersection over Union (IoU)} phạt đồng thời thiếu thông tin
	liên quan và dư thông tin không liên quan:
	\begin{equation}
		\text{IoU}(q;\pi) =
		\frac{|\mathcal{R}(q,k;\pi) \cap \mathcal{G}_q|}
		{|\mathcal{R}(q,k;\pi) \cup \mathcal{G}_q|}
		\label{eq:iou}
	\end{equation}
	IoU là chỉ số tổng hợp quan trọng nhất vì nó yêu cầu đồng thời Recall cao
	\emph{và} Precision cao.
	
	Chỉ số trung bình trên bộ câu hỏi $\mathcal{Q}_{\text{test}}$:
	\begin{equation}
		\overline{\text{IoU}}(\pi) =
		\frac{1}{|\mathcal{Q}_{\text{test}}|}
		\sum_{q \in \mathcal{Q}_{\text{test}}} \text{IoU}(q;\pi)
		\label{eq:mean_iou}
	\end{equation}
	
	\subsection{Chỉ số Tầng Tạo sinh}
	\label{subsec:gen_metrics}
	
	\textbf{Answer Correctness (AC)} đánh giá câu trả lời cuối cùng có đúng
	so với ground truth hay không, đo bằng Exact Match (EM) cho câu hỏi thực
	thể và BERTScore~\cite{zhang2020bertscore} cho câu hỏi tổng hợp:
	\begin{equation}
		\text{BERTScore}(\hat{a}, a^*) =
		F_1\!\left(\text{Prec}(\hat{a}, a^*),\;
		\text{Rec}(\hat{a}, a^*)\right)
		\label{eq:bertscore}
	\end{equation}
	trong đó Precision và Recall được tính trên cơ sở độ tương đồng cosine
	giữa các token embedding BERT của $\hat{a}$ và $a^*$.
	
	% -------------------------------------------------------------------
	\section{Bộ Dữ liệu và Thiết lập Thực nghiệm}
	\label{sec:dataset}
	
	\subsection{Corpus Tuyển sinh HUSC}
	\label{subsec:corpus}
	
	Corpus sử dụng trong thực nghiệm bao gồm bốn nguồn tài liệu chính thức
	của HUSC thu thập năm 2025, như liệt kê tại Bảng~\ref{tab:corpus_stats}.
	Dữ liệu gốc gồm nhiều định dạng (PDF, DOCX, Markdown, TXT, JSONL) và được
	chuẩn hoá về schema thống nhất hai tầng trước khi truy xuất: \texttt{RawDocument}
	(đầu vào trước chunking) và \texttt{RAGChunk} (đầu ra sau chunking).
	
	\begin{table}[htbp]
		\centering
		\caption{Thống kê corpus tuyển sinh HUSC 2025}
		\label{tab:corpus_stats}
		\begin{tabular}{lccc}
			\toprule
			\textbf{Nguồn tài liệu} & \textbf{Số trang} & \textbf{Token (xấp xỉ)}
			& \textbf{Cấu trúc} \\
			\midrule
			Đề án tuyển sinh 2025 & 96  & 42,000  & Chương, điều, khoản \\
			Quy chế tuyển sinh Bộ GD\&ĐT & 52 & 23,000 & Chương, điều \\
			Trang web tuyển sinh HUSC & -- & 18,000  & HTML, bảng \\
			Thông báo bổ sung 2025 & 14 & 6,000  & Văn bản hành chính \\
			\midrule
			\textbf{Tổng cộng} & \textbf{162} & $\mathbf{\approx 89{,}000}$ & \\
			\bottomrule
		\end{tabular}
	\end{table}
	
	\subsection{Bộ Câu hỏi Thực nghiệm}
	\label{subsec:questions}
	
	Bộ câu hỏi thực nghiệm gồm $|\mathcal{Q}_{\text{test}}| = 50$ câu hỏi,
	phân tầng theo ba loại truy vấn đã phân tích tại Bảng~\ref{tab:husc_query_types}
	với tỷ lệ: 20 câu hỏi thực thể đơn (40\%), 18 câu hỏi đa điều kiện (36\%),
	12 câu hỏi so sánh/tổng hợp (24\%). Mỗi câu hỏi được xây dựng kèm theo
	ground-truth answer $a^*_q$ và gold chunk set $\mathcal{G}_q$ xác định
	tay bởi tác giả khoá luận.
	
	\textbf{Vì đây là bộ câu hỏi do tác giả xây dựng và chưa qua đánh giá
		đồng thuận (inter-annotator agreement), các kết quả ở Mục~\ref{sec:results}
		cần được hiểu như bằng chứng định hướng (indicative evidence) chứ không
		phải kết luận tổng quát hóa được.}
	
	\subsection{Thiết lập Hệ thống}
	\label{subsec:system_setup}
	
	Pipeline tiền xử lý được thiết kế theo hướng \emph{canonical-first} nhằm tách
	rõ dữ liệu thô và dữ liệu đã phân đoạn. Cụ thể, dữ liệu đa định dạng được
	chuẩn hoá về \texttt{RawDocument}, sau đó mới đưa qua bộ chunker để sinh
	\texttt{RAGChunk}. Mọi tệp \texttt{chunked\_*.jsonl} đều được chuẩn hoá lại về
	schema thống nhất trước khi ingest vào vector store.
	
	\begin{equation}
		\text{Raw Files} \rightarrow \texttt{RawDocument} \rightarrow \text{Chunker}
		\rightarrow \texttt{RAGChunk} \rightarrow \text{Vector Store}
		\label{eq:canonical_pipeline}
	\end{equation}
	
	Để bảo toàn cấu trúc tài liệu hành chính/pháp lý tiếng Việt, khoá luận áp dụng
	chiến lược \emph{markdown-first chunking}: nhận diện tiêu đề bậc cấu trúc
	(\texttt{PHẦN}, \texttt{CHƯƠNG}, \texttt{MỤC}, \texttt{Điều}, \texttt{Khoản})
	thành heading Markdown trước khi tách theo cửa sổ token có overlap. Cách làm
	này giúp giảm cắt gãy mệnh đề pháp lý và cải thiện khả năng truy xuất theo
	breadcrumbs ngữ nghĩa.
	
	Tất cả thực nghiệm sử dụng cùng một embedding model (\texttt{bge-m3}
	hỗ trợ tiếng Việt), cùng vector store (FAISS), và cùng LLM generation
	model (\texttt{gpt-4o-mini}) để đảm bảo sự khác biệt giữa các chiến lược
	chunking là nguyên nhân duy nhất gây ra sự khác biệt về chỉ số.
	
	\begin{table}[htbp]
		\centering
		\caption{Cấu hình thực nghiệm thống nhất cho tất cả chiến lược chunking}
		\label{tab:exp_config}
		\begin{tabular}{ll}
			\toprule
			\textbf{Thành phần} & \textbf{Cấu hình} \\
			\midrule
			Embedding model & \texttt{BAAI/bge-m3} (1024-dim, hỗ trợ tiếng Việt) \\
			Vector store    & FAISS (IndexFlatIP) \\
			Retrieval top-$k$ & $k = 5$ \\
			LLM generation  & \texttt{gpt-4o-mini} (temperature = 0) \\
			\midrule
			\textit{Tham số riêng cho từng chiến lược:} & \\
			Fixed-size      & $\ell \in \{200, 400, 600\}$, $s = 0.2\ell$ \\
			Recursive       & $\ell \in \{200, 400, 600\}$, $\Sigma$ = mặc định \\
			Semantic        & $p = 0.95$ (phân vị ngưỡng) \\
			Late Chunking   & \texttt{jina-embeddings-v3} ($W = 8{,}192$) \\
			Proposition     & Propositionizer = Flan-T5-large fine-tuned \\
			RAPTOR          & GPT-4o-mini tóm tắt, GMM phân cụm \\
			\bottomrule
		\end{tabular}
	\end{table}
	
	% -------------------------------------------------------------------
	\section{Kết quả Thực nghiệm}
	\label{sec:results}
	
	Mục này trình bày kết quả thực nghiệm trên bộ câu hỏi tuyển sinh HUSC.
	Bảng~\ref{tab:main_results} tổng hợp ba chỉ số chính ($\overline{\text{Recall}@5}$,
	$\overline{\text{Precision}@5}$, $\overline{\text{IoU}}$) và Answer
	Correctness cho tất cả chiến lược được so sánh.
	
	\begin{table}[htbp]
		\centering
		\caption{Kết quả thực nghiệm so sánh các chiến lược chunking trên bộ câu
			hỏi tuyển sinh HUSC ($n = 50$ câu hỏi, $k = 5$, LLM = gpt-4o-mini).
			\textbf{In đậm} = kết quả tốt nhất trên từng chỉ số.
			\dag{} = kết quả của tác giả; các số liệu khác cần được điền sau khi
			hoàn thành thực nghiệm.}
		\label{tab:main_results}
		\begin{adjustbox}{max width=\linewidth}
			\begin{tabular}{lcccc}
				\toprule
				\textbf{Chiến lược}
				& $\overline{\text{Recall}@5}$
				& $\overline{\text{Precision}@5}$
				& $\overline{\text{IoU}}$
				& \textbf{AC (BERTScore)} \\
				\midrule
				Fixed-size ($\ell$=200)  & -- & -- & -- & -- \\
				Fixed-size ($\ell$=400)  & -- & -- & -- & -- \\
				Fixed-size ($\ell$=600)  & -- & -- & -- & -- \\
				Recursive ($\ell$=400)   & -- & -- & -- & -- \\
				Semantic (baseline)      & -- & -- & -- & -- \\
				Late Chunking            & -- & -- & -- & -- \\
				Proposition-based        & -- & -- & -- & -- \\
				RAPTOR (3 levels)        & -- & -- & -- & -- \\
				\midrule
				\textbf{Hybrid} (Semantic + Proposition) & -- & -- & -- & -- \\
				\bottomrule
			\end{tabular}
		\end{adjustbox}
	\end{table}
	
	\noindent\textit{Ghi chú: Các ô ``--'' sẽ được điền bằng kết quả thực nghiệm
		của tác giả khi hoàn thành triển khai. Không có số liệu nào được đặt trước
		trong khoá luận này; mọi con số trong bảng kết quả sẽ xuất phát từ chạy
		thực nghiệm trực tiếp trên corpus HUSC với hệ thống được mô tả tại
		Bảng~\ref{tab:exp_config}.}
	
	% -------------------------------------------------------------------
	\section{Phân tích Kết quả theo Loại Truy vấn}
	\label{sec:analysis_by_type}
	
	Ngoài chỉ số tổng hợp, khoá luận phân tích kết quả tách biệt theo ba loại
	truy vấn đã định nghĩa tại Mục~\ref{sec:husc_problem}. Phân tích này quan
	trọng vì một chiến lược có thể đạt IoU trung bình cao nhờ vượt trội ở loại
	truy vấn thực thể đơn (40\% tập câu hỏi) trong khi thất bại ở loại câu hỏi
	so sánh (25\%).
	
	Dựa trên lý thuyết tại Chương~\ref{chap:chunking_theory}, các giả thuyết
	thực nghiệm cụ thể được phát biểu như sau:
	
	\begin{menh}[Giả thuyết H1]
		Với câu hỏi thực thể đơn, Proposition-based chunking đạt Recall@5 cao hơn
		Fixed-size ($\ell = 400$) ít nhất $5\%$ (tuyệt đối), phù hợp với kết quả
		$+10.1\%$ trên unsupervised retrievers của Chen et al.~\cite{chen2024dense}
		--- tuy nhiên không thể kỳ vọng cùng biên độ vì retriever trong thực
		nghiệm này (\texttt{bge-m3}) là supervised model và corpus tiếng Việt khác
		về phân phối so với Wikipedia tiếng Anh.
	\end{menh}
	
	\begin{menh}[Giả thuyết H2]
		Với câu hỏi so sánh/tổng hợp, RAPTOR đạt Answer Correctness (BERTScore)
		cao hơn tất cả chiến lược G1--G3, nhờ cơ chế tree collapse cho phép truy
		xuất đồng thời ở nhiều cấp trừu tượng.
	\end{menh}
	
	\begin{menh}[Giả thuyết H3]
		Chiến lược Hybrid (Semantic chunking cho vector search kết hợp Proposition
		extraction cho đồ thị tri thức) đạt $\overline{\text{IoU}}$ tổng hợp cao
		nhất, đánh đổi giữa chi phí tính toán và chất lượng trên cả ba loại truy
		vấn.
	\end{menh}
	
	% -------------------------------------------------------------------
	\section{Phân tích Chi phí -- Lợi ích}
	\label{sec:cost_benefit}
	
	Hình~\ref{fig:latency_quality} trình bày không gian đánh đổi giữa
	$\overline{\text{IoU}}$ (chất lượng truy xuất) và thời gian xây dựng chỉ
	mục (indexing latency) cho mỗi chiến lược. Đây là thước đo thực tiễn quan
	trọng: trong môi trường học thuật với tài nguyên GPU hạn chế, một chiến
	lược đạt IoU cao hơn 5\% nhưng tốn thời gian index gấp 10 lần có thể không
	khả thi cho triển khai thực tế.
	
	\begin{figure}[htbp]
		\centering
		\begin{adjustbox}{max width=\linewidth}
			\begin{tikzpicture}[
				row1/.style={rectangle, draw=black!30, fill=gray!8,
					minimum height=0.75cm, align=center, font=\footnotesize},
				rowG/.style={rectangle, draw=black!30,
					minimum height=0.75cm, align=center, font=\footnotesize},
				hdr/.style={rectangle, draw=black!50, fill=black!12,
					minimum height=0.75cm, align=center, font=\footnotesize\bfseries},
				bar/.style={rectangle, draw=none, minimum height=0.55cm},
				]
				% Header
				\node[hdr, text width=2.8cm] (h1) at (0,    0) {Chiến lược};
				\node[hdr, text width=3.5cm] (h2) at (3.3,  0) {Chi phí indexing};
				\node[hdr, text width=4.0cm] (h3) at (7.25, 0) {IoU (sau thực nghiệm)};
				
				% Rows
				\foreach \i/\strat/\cost/\costw/\col in {
					1/Fixed-size (G1)/Rất thấp/0.5cm/gray!50,
					2/Recursive (G1)/Thấp/1.0cm/gray!60,
					3/Semantic (G2)/Trung bình/1.8cm/blue!50,
					4/Late Chunking (G3)/Trung bình/1.8cm/orange!60,
					5/Proposition (G3)/Cao/2.8cm/orange!70,
					6/RAPTOR (G4)/Cao/2.8cm/red!55,
					7/Contextual (G4)/Rất cao/3.8cm/red!70
				}{
					\node[rowG, text width=2.8cm, fill=white,
					draw=black!20] at (0, -\i*0.8)
					{\strat};
					\node[rowG, text width=3.5cm, fill=white,
					draw=black!20] at (3.3, -\i*0.8)
					{\cost};
					% Bar placeholder for IoU
					\node[rowG, text width=4.0cm, fill=gray!5,
					draw=black!20] at (7.25, -\i*0.8)
					{\textit{--}};
					% Cost bar
					\node[bar, minimum width=\costw, fill=\col, opacity=0.7]
					at ({1.6 + \costw/2 - 0.1cm}, -\i*0.8) {};
				}
				
				% Footnote line
				\node[font=\tiny\itshape, text=gray!70, text width=9cm]
				at (3.6, -6.5)
				{Cột IoU sẽ được điền bằng kết quả thực nghiệm.
					Chi phí indexing là thứ bậc định tính, không phải số tuyệt đối.};
			\end{tikzpicture}
		\end{adjustbox}
		\caption{Đánh đổi chi phí indexing -- chất lượng truy xuất ($\overline{\text{IoU}}$)
			của các chiến lược chunking. Cột IoU sẽ được điền bằng kết quả thực
			nghiệm; cột chi phí indexing là thứ bậc định tính dựa trên phân tích
			độ phức tạp tính toán lý thuyết.}
		\label{fig:latency_quality}
	\end{figure}
	
	% -------------------------------------------------------------------
	\section{Hạn chế Dữ liệu và Tiền xử lý}
	\label{sec:data_preprocess_limits}
	
	Mặc dù pipeline chuẩn hoá đã giảm đáng kể sai khác cấu trúc giữa các nguồn,
	một số hạn chế kỹ thuật vẫn tồn tại. Thứ nhất, trích xuất văn bản từ PDF dựa
	trên lớp text có thể suy giảm chất lượng với tài liệu scan ảnh; trong trường
	hợp này, chunking chịu lỗi lan truyền từ bước OCR. Thứ hai, tài liệu DOCX có
	thể chứa thông tin quan trọng trong bảng, header/footer hoặc textbox; các thành
	phần này không phải lúc nào cũng được bảo toàn đầy đủ khi chuyển sang text phẳng.
	
	Với văn bản pháp lý/hành chính, khoá luận áp dụng làm sạch heuristic (loại
	header/footer lặp, nhận diện cấu trúc \texttt{Điều/Khoản/Điểm}). Tuy nhiên,
	heuristic vẫn có rủi ro loại nhầm trong các mẫu trình bày không chuẩn hoá.
	Do đó, kết quả thực nghiệm cần được hiểu trong bối cảnh pipeline hiện tại, và
	một hướng mở quan trọng là tái tạo corpus từ bản gốc có cấu trúc giàu hơn
	(HTML/PDF có layout) thay vì chỉ dựa trên text extraction tuyến tính.
	
	% -------------------------------------------------------------------
	\section{Xu hướng Chunking năm 2026 và Hàm ý cho HUSC}
	\label{sec:trend_2026_chunking}
	
	Các kết quả gần đây cho thấy không còn xu hướng chọn một kỹ thuật chunking duy
	nhất cho toàn hệ thống. Thay vào đó, hướng đi nổi bật là kiến trúc đa tầng:
	chunk nhỏ để tăng độ nhạy truy xuất, đồng thời duy trì chunk cha lớn hơn để
	khôi phục ngữ cảnh khi cần. Phân tích của Merola \& Singh cho thấy các kỹ thuật
	``giữ ngữ cảnh'' như late/contextual chunking cải thiện rõ rệt chất lượng trả
	lời nhưng đổi lại là chi phí indexing cao~\cite{merola2025reconstructing}.
	Ở một hướng khác, HiChunk chỉ ra rằng biểu diễn phân cấp giúp hệ thống cân bằng
	được giữa câu hỏi factoid và câu hỏi tổng hợp dài hơi~\cite{lu2025hichunk}.
	
	Trên bình diện triển khai, các framework RAG năm 2026 đều hội tụ về cùng một
	mẫu thiết kế: tạo cây chunk nhiều mức (coarse-to-fine), sau đó hợp nhất ngữ
	cảnh theo truy vấn thay vì cố định ngay từ đầu~\cite{llamaindex2026docs}.
	Với bộ tách recursive, thực hành phổ biến vẫn là giữ overlap vừa phải để giảm
	hiệu ứng biên, rồi nâng cấp bằng semantic/hierarchical khi bài toán đòi hỏi
	tính suy luận cao hơn~\cite{langchain2026docs}.
	
	% -------------------------------------------------------------------
	\section{Đề xuất Chiến lược cho Bài toán HUSC}
	\label{sec:recommendation}
	
	Từ đặc thù dữ liệu tuyển sinh (nhiều điều kiện, nhiều ngoại lệ, nhiều văn bản
	pháp lý), khoá luận không chọn một ``siêu chiến lược'' duy nhất. Thay vào đó,
	hệ thống dùng chunking thích ứng theo truy vấn, trên cùng một corpus đã chuẩn
	hoá markdown-first.
	
	\textbf{Lớp 1 --- Retrieval hạt mịn.} Tài liệu được tách semantic/recursive ở
	mức $\ell \approx 200$--$420$ token (tuỳ profile), overlap khoảng $10$--$25\%$.
	Lớp này tối ưu cho câu hỏi thực thể và tra cứu điều kiện ngắn.
	
	\textbf{Lớp 2 --- Retrieval ngữ cảnh phân cấp.} Trên cùng tài liệu, hệ thống
	duy trì cây chunk cha--con (ví dụ các mức 2048$\rightarrow$512$\rightarrow$128)
	để có thể ``nâng ngữ cảnh'' khi truy vấn cần đọc xuyên nhiều đoạn~\cite{llamaindex2026docs}.
	Điều kiện hợp nhất nút cha được mô tả bởi:
	\begin{equation}
		r(p)=\frac{|\mathcal{C}_{\text{hit}}(p)|}{|\mathcal{C}(p)|},
		\qquad
		r(p)>\tau_{\text{merge}} \Rightarrow p\ \text{được đưa vào ngữ cảnh}
		\label{eq:auto_merge_ratio}
	\end{equation}
	trong đó $\tau_{\text{merge}}$ là ngưỡng hợp nhất (thực nghiệm khởi tạo ở mức
	$0.5$).
	
	\textbf{Lớp 3 --- Proposition/Graph cho truy vấn đa bước.} Với câu hỏi so sánh
	hoặc cần liên kết nhiều fact rời, hệ thống dùng thêm lớp proposition để tạo
	đồ thị tri thức, sau đó hợp nhất điểm xếp hạng với dense retrieval bằng RRF.
	
	Về điều phối truy vấn, hệ thống dùng trọng số động theo loại câu hỏi:
	\begin{equation}
		\mathbf{w}(q) = \big(w_f(q),\, w_h(q),\, w_p(q)\big),
		\quad
		w_f+w_h+w_p=1
		\label{eq:adaptive_weight}
	\end{equation}
	trong đó $w_f$ ưu tiên lớp hạt mịn, $w_h$ ưu tiên lớp phân cấp, và $w_p$ ưu
	tiên lớp proposition. Với truy vấn thực thể, $w_f$ thường lớn nhất; với truy
	vấn so sánh/tổng hợp, $w_h$ và $w_p$ được tăng lên để tránh thiếu ngữ cảnh.
	
	Lựa chọn này phù hợp với xu hướng 2026: \emph{hybrid + hierarchical + adaptive}
	thay vì cố ép toàn bộ hệ thống vào một kích thước chunk cố định.
	
	\clearpage
	\addcontentsline{toc}{chapter}{Tài liệu tham khảo}
	
	\begin{thebibliography}{99}
		
		\bibitem{vaswani2017attention}
		Vaswani, A., Shazeer, N., Parmar, N., Uszkoreit, J., Jones, L., Gomez, A.~N.,
		Kaiser, L., \& Polosukhin, I. (2017).
		\textit{Attention Is All You Need}.
		\textit{Advances in Neural Information Processing Systems} (NeurIPS), 30,
		5998--6008.
		
		\bibitem{lewis2020retrieval}
		Lewis, P., Perez, E., Piktus, A., Petroni, F., Karpukhin, V., Goyal, N.,
		Küttler, H., Lewis, M., Yih, W., Rocktäschel, T., Riedel, S., \& Kiela, D.
		(2020).
		\textit{Retrieval-Augmented Generation for Knowledge-Intensive NLP Tasks}.
		\textit{Advances in Neural Information Processing Systems} (NeurIPS), 33,
		9459--9474.
		
		\bibitem{barnett2024seven}
		Barnett, S., Kurniawan, S., Thudumu, S., Brannelly, Z., \& Abdelrazek, M.
		(2024).
		\textit{Seven Failure Points When Engineering a Retrieval Augmented Generation
			System}.
		In \textit{Proceedings of the 3rd International Workshop on Software
			Engineering and AI for Data Quality (SEA4DQ)}, co-located with ICSE 2024.
		arXiv:2401.05856.
		
		\bibitem{chroma2024eval}
		Chroma Research. (2024).
		\textit{Evaluating Chunking Strategies for Retrieval}.
		Technical Report.
		\url{https://research.trychroma.com/evaluating-chunking}
		
		\bibitem{gunther2024late}
		Günther, M., Mohr, I., Williams, D.~J., Wang, B., \& Xiao, H. (2024).
		\textit{Late Chunking: Contextual Chunk Embeddings Using Long-Context Embedding
			Models}.
		arXiv preprint arXiv:2409.04701.
		
		\bibitem{chen2024dense}
		Chen, T., Wang, H., Chen, S., Yu, W., Ma, K., Zhao, X., Zhang, H., \& Yu, D.
		(2024).
		\textit{Dense X Retrieval: What Retrieval Granularity Should We Use?}
		In \textit{Proceedings of the 2024 Conference on Empirical Methods in Natural
			Language Processing} (EMNLP), Miami, Florida, pp.\ 15159--15177.
		Association for Computational Linguistics.
		\url{https://aclanthology.org/2024.emnlp-main.845/}
		
		\bibitem{sarthi2024raptor}
		Sarthi, P., Abdullah, S., Tuli, A., Khanna, S., Goldie, A., \& Manning,
		C.~D. (2024).
		\textit{RAPTOR: Recursive Abstractive Processing for Tree-Organized Retrieval}.
		In \textit{Proceedings of the 12th International Conference on Learning
			Representations} (ICLR 2024).
		\url{https://openreview.net/forum?id=GN921JHCRw}
		
		\bibitem{anthropic2024contextual}
		Anthropic. (2024).
		\textit{Contextual Retrieval}.
		Technical Blog.
		\url{https://www.anthropic.com/news/contextual-retrieval}
		
		\bibitem{merola2025reconstructing}
		Merola, C., \& Singh, J. (2025).
		\textit{Reconstructing Context: Evaluating Advanced Chunking Strategies for
			Retrieval-Augmented Generation}.
		arXiv preprint arXiv:2504.19754.
		Presented at the Second Workshop on Knowledge-Enhanced Information Retrieval,
		ECIR 2025.
		
		\bibitem{lu2025hichunk}
		Lu, W., Chen, K., Qiao, R., \& Sun, X. (2025).
		\textit{HiChunk: Evaluating and Enhancing Retrieval-Augmented Generation
			with Hierarchical Chunking}.
		arXiv preprint arXiv:2509.11552.
		
		\bibitem{llamaindex2026docs}
		LlamaIndex. (2026).
		\textit{Hierarchical Node Parser and Auto-Merging Retriever Documentation}.
		Official Documentation.
		\url{https://developers.llamaindex.ai/python/framework/module_guides/loading/node_parsers/modules}
		
		\bibitem{langchain2026docs}
		LangChain. (2026).
		\textit{RecursiveCharacterTextSplitter Documentation}.
		Official Documentation.
		\url{https://docs.langchain.com/oss/python/integrations/splitters/recursive_text_splitter}
		
		\bibitem{robertson2009bm25}
		Robertson, S., \& Zaragoza, H. (2009).
		\textit{The Probabilistic Relevance Framework: BM25 and Beyond}.
		\textit{Foundations and Trends in Information Retrieval}, 3(4), 333--389.
		
		\bibitem{cormack2009rrf}
		Cormack, G.~V., Clarke, C.~L.~A., \& Buettcher, S. (2009).
		\textit{Reciprocal Rank Fusion Outperforms Condorcet and Individual Rank
			Learning Methods}.
		In \textit{Proceedings of the 32nd Annual International ACM SIGIR Conference},
		pp.\ 758--759.
		
		\bibitem{zhang2020bertscore}
		Zhang, T., Kishore, V., Wu, F., Weinberger, K.~Q., \& Artzi, Y. (2020).
		\textit{BERTScore: Evaluating Text Generation with BERT}.
		In \textit{Proceedings of the 8th International Conference on Learning
			Representations} (ICLR 2020).
		
		\bibitem{ethayarajh2019contextual}
		Ethayarajh, K. (2019).
		\textit{How Contextual are Contextualized Word Representations?
			Comparing the Geometry of BERT, ELMo, and GPT-2 Embeddings}.
		In \textit{Proceedings of the 2019 Conference on Empirical Methods in
			Natural Language Processing} (EMNLP-IJCNLP), pp.\ 55--65.
		
	\end{thebibliography}
	
\end{document}